Goal: test H1 with two nested logistic models that differ only by the block $\mathbf{Hd}$.

\section{Pipeline}
\begin{enumerate}
    \item Pre-specified protocol: outcome, regions, features, thresholds, folds, seeds.
    \item Cohort and outcome.
    \item Hemodynamic features $\mathbf{Hd}$.
    \item Pre-validated benchmark score $S_B$.
    \item Inferential test (LRT) on the full cohort.
    \item Out-of-fold performance estimates.
    \item Sensitivity analyses.
\end{enumerate}

\section{Cohort and outcome}
\begin{itemize}
    \item Patients treated with EVAR, with a valid preoperative CFD (Appendix~A) and benchmark features.
    \item Expected size: about 100--150 patients. [Open point: number of events.]
    \item Minimum follow-up (e.g.\ 12 months) is the only outcome-related exclusion: it avoids false negatives.
    \item Patients without complications are kept; imbalance is handled by the model intercept.
    \item One binary primary outcome $y_i \in \{0,1\}$, fixed in advance: any endoleak, or any complication.
    \item Subtypes (type I, type II) only as exploratory analyses.
\end{itemize}

\section{Hemodynamic features}
\begin{itemize}
    \item Point indices TAWSS, OSI, RRT on the last cycle (Section~\ref{sec:wss-indices}).
    \item Anatomical regions $R_k$ from the centerline: proximal neck, sac, iliac landing zones.
\end{itemize}
\begin{equation}
    \bar{q}_{k} = \frac{1}{|R_k|}\int_{R_k} q(\mathbf{x})\,dA,
    \qquad
    f_{k}(\theta) = \frac{1}{|R_k|}\int_{R_k} \mathbb{1}\left[q(\mathbf{x}) > \theta\right] dA
\end{equation}
\begin{itemize}
    \item $\mathbf{Hd}_i \in \mathbb{R}^p$, with $p = 2$--$5$, chosen in advance from a mechanistic hypothesis, not from the outcome.
    \item Thresholds $\theta$ from the literature. \citeneeded{TAWSS and OSI thresholds}
    \item Features standardised before fitting.
\end{itemize}

\section{Benchmark score}
\begin{itemize}
    \item Benchmark: the supervisor's SVM on radiomics and morphology.
    \item Summarised in one number per patient: one parameter instead of many correlated features.
    \item Pre-validation: 5 stratified folds, $R = 10$--$20$ seeds; the whole benchmark pipeline (feature selection, tuning, SVM) is refitted inside each fold. \citeneeded{Tibshirani and Efron 2002}
    \item If the benchmark was developed on another cohort, its predictions are used directly.
\end{itemize}
\begin{equation}
    \bar{p}_i = \frac{1}{R}\sum_{r=1}^{R} \hat{p}_i^{(r)},
    \qquad
    S_{B,i} = \operatorname{logit}(\bar{p}_i) = \log\frac{\bar{p}_i}{1-\bar{p}_i}
\end{equation}

\section{Models}
\begin{align}
    M_0:\;& \operatorname{logit} P(y_i = 1) = \beta_0 + \beta_1 S_{B,i} \\
    M_1:\;& \operatorname{logit} P(y_i = 1) = \beta_0 + \beta_1 S_{B,i} + \boldsymbol{\gamma}^{\top}\mathbf{Hd}_i
\end{align}
\begin{itemize}
    \item $M_0$ is the logistic recalibration of the benchmark: same ranking, same AUC, better calibrated. \citeneeded{Cox 1958}
    \item $M_1$ is a model extension \cite{janssen_updating_2008}; $\beta_1$ is re-estimated together with $\boldsymbol{\gamma}$.
\end{itemize}

\section{Hypothesis and test}
\begin{equation}
    H_0^{(1)}:\; \boldsymbol{\gamma} = \mathbf{0}
    \qquad \text{vs.} \qquad
    H_A^{(1)}:\; \boldsymbol{\gamma} \neq \mathbf{0}
\end{equation}
\begin{equation}
    \Lambda_1 = 2\left(\ell_1 - \ell_0\right) \;\overset{H_0}{\sim}\; \chi^2_{p}
\end{equation}
\begin{itemize}
    \item $\ell_j$: maximised log-likelihood of $M_j$, fitted on the same patients.
    \item Asymptotic $\chi^2_p$ by Wilks' theorem. \citeneeded{Wilks 1938}
    \item $\alpha = 0.05$; Holm--Bonferroni if H1 and H2 are both primary.
    \item Choice of the LRT: Chapter~\ref{ch:framework} \cite{pepe_testing_2013}.
\end{itemize}

\section{Performance}
Out-of-fold chain: $M_0$ and $M_1$ are refitted inside the same folds, giving predictions $\hat{\pi}^{(0)}_i$ and $\hat{\pi}^{(1)}_i$.
\begin{equation}
    \Delta\mathrm{AUC} = \mathrm{AUC}\left(\hat{\pi}^{(1)}\right) - \mathrm{AUC}\left(\hat{\pi}^{(0)}\right)
\end{equation}
\begin{equation}
    \Delta\mathrm{BS} = \frac{1}{n}\sum_{i=1}^{n}\left[\left(\hat{\pi}_i^{(1)} - y_i\right)^2 - \left(\hat{\pi}_i^{(0)} - y_i\right)^2\right]
\end{equation}
\begin{equation}
    \mathrm{NB}(p_t) = \frac{\mathrm{TP}(p_t)}{n} - \frac{\mathrm{FP}(p_t)}{n}\cdot\frac{p_t}{1 - p_t}
    \quad \citeneeded{Vickers and Elkin 2006}
\end{equation}
\begin{itemize}
    \item Calibration intercept and slope out of fold \cite{alba_discrimination_2017}.
    \item 95\% CIs by paired bootstrap ($B = 2000$).
    \item Effect sizes, not tests.
\end{itemize}

\section{Sample size}
\begin{equation}
    1 + p \;\le\; \frac{E}{10}
    \quad \citeneeded{Peduzzi et al.\ 1996}
\end{equation}
\begin{itemize}
    \item $E$: number of events. With $E = 40$, at most $p = 3$.
    \item The 10-events rule is a heuristic. \citeneeded{Riley et al.\ 2020}
\end{itemize}

\section{Sensitivity analyses}
\begin{itemize}
    \item Spline on $S_B$: checks linearity in the logit.
    \item Firth regression: small-sample bias. \citeneeded{Firth 1993}
    \item Parametric bootstrap of the LRT: checks the $\chi^2$ approximation. \citeneeded{Höfling and Tibshirani 2008}
    \item Wald test, number of seeds, follow-up threshold.
    \item Exploratory: $S_B \times \mathbf{Hd}$ interaction, other outcomes.
\end{itemize}

\section{Limitations}
\label{sec:p1-limits}
\begin{itemize}
    \item $S_B$ is assumed to carry all useful benchmark information; interactions with single features are lost.
    \item Pre-validation is not a truly independent set.
    \item $S_B$ is estimated but treated as fixed: p-values slightly optimistic.
    \item Few events: low power; asymptotic approximations less reliable.
    \item Binary outcome with minimum follow-up: time-to-event and competing risks ignored.
    \item Selection bias from excluded patients (failed CFD, short follow-up).
    \item CFD measurement error (boundary conditions, rigid walls) biases $\boldsymbol{\gamma}$ towards zero.
    \item Collinearity between $\mathbf{Hd}$ and morphology.
    \item Bootstrap on out-of-fold predictions ignores refitting variability.
    \item Internal validation only, single centre; association, not causation.
\end{itemize}
